\clearpage
\fancyhead[L]{\small VESTA. CERES. COMET.}
\SourceLeaf{313}\begin{center}
VESTA. CERES. COMET.
\end{center}

\par\vspace{4pt}\hrule\vspace{4pt}
\begin{center}
\fontsize{10}{12}\selectfont Göttingische gelehrte Anzeigen. Issue 180. Pp. 1793..1795. 10 November 1808.
\end{center}
\par\vspace{4pt}\hrule\vspace{4pt}

Professor \textsc{Bode} has sent us the following valuable observations of \emph{Vesta}, made at the Berlin Observatory:


\begin{center}\fontsize{10}{12}\selectfont\setlength{\tabcolsep}{5pt}\renewcommand{\arraystretch}{1.15}\begin{tabular}{lrrr}
1808 & \shortstack{Mean time\\at Berlin} & \shortstack{Apparent right\\ascension} & \shortstack{Apparent south\\declination}\\
\hline
Aug. 24 & $11^{\text{u}}\;16^{\text{m}}\;16^{\text{s}}$ & $354^\circ\;34'\;59''$ & $13^\circ\;50'\;22''$\\
Sept. 1 & $12\;48\;0.7$ & $352\;57\;27$ & $14\;49\;41$\\
5 & $11\;8\;0$ & $352\;5\;20$ & $15\;20\;11$\\
13 & $11\;50\;1.0$ & $350\;15\;10$ & $16\;15\;59$\\
15 & $11\;40\;21.7$ & $349\;48\;2$ & $16\;28\;2$\\
20 & $11\;16\;26.7$ & $348\;42\;5$ & $16\;54\;35$\\
Oct. 6 & $10\;1\;55.2$ & $345\;49\;31$ & $17\;42\;7$\\
9 & $9\;48\;36.0$ & $345\;26\;47$ & $17\;44\;4$\\
14 & $9\;26\;57.3$ & $344\;56\;58$ & $17\;42\;30$\\
\end{tabular}\end{center}


\par\indent All these observations except the first and third were made with the mural quadrant.

\par\indent \emph{Ceres} remained invisible with the mural quadrant throughout. The following three observations were made with the ring micrometer:


\begin{center}\fontsize{10}{12}\selectfont\setlength{\tabcolsep}{5pt}\renewcommand{\arraystretch}{1.15}\begin{tabular}{lrrr}
1808 & \shortstack{Mean time\\at Berlin} & \shortstack{Apparent right\\ascension} & \shortstack{Apparent south\\declination}\\
\hline
Aug. 29 & $10^{\text{u}}\;52^{\text{m}}\;49^{\text{s}}$ & $314^\circ\;45'\;57''$ & $30^\circ\;34'\;28''$\\
Sept. 16 & $7\;46\;36$ & $312\;26\;31$ & $30\;39\;58$\\
17 & $9\;12\;43$ & $312\;22\;24$ & $30\;40\;4$\\
\end{tabular}\end{center}


\par\indent The following observations of the most recent comet were made in Petersburg by Messrs. \textsc{Schubert} and \textsc{Wisniewsky} and were also kindly sent to us by Professor \textsc{Bode}. They are of particular interest because they are the latest that could be made anywhere. By the time of the final observations, the comet was so faint that only a very sharp eye could make it out in the telescope. The observations were made with the ring micrometer. The reduction took no account of refraction, nutation or aberration. The comparison stars were not specified.


\begin{center}\fontsize{10}{12}\selectfont\setlength{\tabcolsep}{5pt}\renewcommand{\arraystretch}{1.15}\begin{tabular}{lrrr}
1808 & \shortstack{Mean time\\at Petersburg} & \shortstack{Right\\ascension} & \shortstack{North\\declination}\\
\hline
Jan. 15 & $6^{\text{u}}\;49^{\text{m}}\;30^{\text{s}}$ & $335^\circ\;50'\;2''$ & $47^\circ\;47'\;27''$\\
26 & $7\;2\;25$ & $345\;33\;24$ & $48\;4\;6$\\
Feb. 15 & $8\;45\;44$ & $2\;36\;58$ & $48\;44\;20$\\
25 & $9\;17\;3$ & $9\;21\;23$ & $48\;50\;48$\\
Mar. 1 & $9\;53\;5$ & $12\;33\;34$ & $48\;51\;59$\\
22 & $9\;24\;29$ & $23\;6\;39$ & $48\;48\;25$\\
27 & $10\;54\;20$ & $25\;48\;51$ & $48\;53\;49$\\
\end{tabular}\end{center}


\par\indent We also take this opportunity to include the observations of \emph{Vesta} made at our observatory this year. The first, which appears to be the earliest ma\SourceLeaf{314}de anywhere that year, has already been reported in issue 107 of this journal [p. 303 of this volume]. From 10 August onward the observations were made with the mural quadrant. Unforeseen obstacles interrupted the series in October.


\begin{center}\fontsize{10}{12}\selectfont\setlength{\tabcolsep}{5pt}\renewcommand{\arraystretch}{1.15}\begin{tabular}{lrrr}
1808 & \shortstack{Mean time\\at Göttingen} & \shortstack{Apparent right\\ascension} & \shortstack{Apparent south\\declination}\\
\hline
June 22 & $13^{\text{u}}\;46^{\text{m}}\;25^{\text{s}}$ & $353^\circ\;12'\;17''$ & $9^\circ\;12'\;16''$\\
July  1 & $14\;38\;57$ & $354\;28\;3$ & $9\;9\;39$\\
2 & $12\;50\;2$ & $355\;7\;7$ & \\
13 & $13\;22\;37$ & $356\;35\;27$ & $9\;27\;40$\\
Aug. 10 & $14\;28\;53.7$ & $356\;36\;16.5$ & $11\;52\;50.3$\\
12 & $14\;20\;22.3$ & $356\;23\;22.3$ & $12\;8\;5.3$\\
14 & $14\;11\;23.4$ & $356\;9\;3.7$ & $12\;23\;25.7$\\
15 & $14\;6\;54.8$ & $356\;1\;6.0$ & $12\;31\;21.2$\\
22 & $13\;35\;4.7$ & $354\;56\;6.3$ & $13\;28\;20.7$\\
23 & $13\;30\;28.0$ & $354\;45\;24.1$ & $13\;36\;10.4$\\
29 & $13\;2\;11.5$ & $353\;35\;29.2$ & $14\;25\;25.8$\\
30 & $12\;57\;26.4$ & $353\;23\;12.9$ & $14\;33\;32.6$\\
Sept. 5 & $12\;28\;37.0$ & $352\;4\;33.3$ & $15\;20\;32.5$\\
8 & $12\;14\;6.6$ & $351\;23\;46.5$ & $15\;42\;22.3$\\
11 & $11\;59\;34.6$ & $350\;43\;7.3$ & $16\;3\;7.1$\\
12 & $11\;54\;42.9$ & $350\;28\;40.5$ & $16\;9\;34.9$\\
21 & $11\;11\;22.7$ & $348\;29\;10.5$ & $16\;59\;51.1$\\
\end{tabular}\end{center}


\Needspace{17\baselineskip}
\par\indent Mr. \textsc{Tiarks}, who has been pursuing the mathematical sciences here with success and already has considerable skill in astronomical calculation, compared some of these observations with the elements we published in issue 107 of these Anzeigen [p. 303 of this volume]. The comparison gave the following differences:


\begin{center}\fontsize{10}{12}\selectfont\setlength{\tabcolsep}{5pt}\renewcommand{\arraystretch}{1.15}\begin{tabular}{lrr}
 & \multicolumn{2}{c}{Difference}\\
 & in longitude & in latitude\\\hline
Aug. 23 & $+4'\;7''4$ & $+43''0$\\
29 & $+4\;5.6$ & $+50.6$\\
30 & $+3\;53.9$ & $+45.0$\\
Sept. 5 & $+4\;4.5$ & $+45.2$\\
8 & $+4\;11.3$ & $+54.0$\\
11 & $+3\;31.6$ & $+38.3$\\
12 & $+4\;7.7$ & $+59.6$\\
21 & $+3\;37.0$ & $+55.1$\\
\end{tabular}\end{center}


